%───────────────────────────────────────────────────────────────────────────────
\chapter{Figuras y gráficas}
\label{ch:figuras}
%───────────────────────────────────────────────────────────────────────────────

Las figuras en un documento académico cumplen dos funciones distintas:
las gráficas generadas por \texttt{pgfplots}~\cite{pgfplots2023} son parte
del argumento matemático --- muestran convergencia, distribuciones, campos
de fase --- y deben ser tan rigurosas como una ecuación. Las imágenes
externas importadas con \texttt{graphicx}~\cite{graphicx2021} son apoyo
visual: diagramas de arquitectura, fotografías de montajes, resultados de
simulaciones externas. El tratamiento tipográfico es el mismo pero su
naturaleza es diferente.

%───────────────────────────────────────────────────────────────────────────────
\section{Figuras externas}
\label{sec:figuras-externas}
%───────────────────────────────────────────────────────────────────────────────

\subsection{Figura simple}
\label{subsec:figura-simple}

La estructura canónica de una figura flotante en flujo normal de texto:

\begin{lstlisting}[style=latex, frame=none, numbers=none]
\begin{figure}[htbp]
  \centering
  \includegraphics[width=0.8\textwidth]{figs/nombre.pdf}
  \caption{Descripcion completa que se lee sin ver la figura.}
  \label{fig:etiqueta}
\end{figure}
\end{lstlisting}

\begin{observacion}{Formato de imagen recomendado}
LuaLaTeX acepta PDF, PNG y JPEG vía \texttt{\textbackslash includegraphics}.
Para gráficas vectoriales (diagramas, plots generados externamente), siempre
usar PDF — escala sin pérdida de calidad a cualquier resolución de impresión.
Para fotografías o capturas de pantalla, PNG. Los SVG deben convertirse a
PDF primero (\texttt{inkscape -{}-export-filename=out.pdf in.svg}).
\end{observacion}

\subsection{Figura dentro de tcolorbox}
\label{subsec:figura-en-caja}

Dentro de cualquier entorno basado en \texttt{tcolorbox}, usar \texttt{boxfig}
en lugar de \texttt{figure}. El ancho de referencia cambia de
\texttt{\textbackslash textwidth} a \texttt{\textbackslash linewidth}:

\begin{lstlisting}[style=latex, frame=none, numbers=none]
\begin{ejemplo}{Campo electrico de una carga puntual}
  \begin{boxfig}{fig:campo-e}{Lineas de campo de $\eqvec{\mathbf{E}}$}
    \includegraphics[width=0.7\linewidth]{figs/campo_electrico.pdf}
  \end{boxfig}
  La intensidad decae como $\eqfn{E}(r) \propto \eqk{q}/\equ{r}^2$.
\end{ejemplo}
\end{lstlisting}

\subsection{Subfiguras}
\label{subsec:subfiguras}

El paquete \texttt{subcaption}~\cite{subcaption2022} permite agrupar
múltiples figuras con subcaptions independientes:

\begin{lstlisting}[style=latex, frame=none, numbers=none]
\begin{figure}[htbp]
  \centering
  \begin{subfigure}[b]{0.48\textwidth}
    \centering
    \includegraphics[width=\textwidth]{figs/solucion_t0.pdf}
    \caption{$t = 0$: condicion inicial.}
    \label{fig:calor-t0}
  \end{subfigure}
  \hfill
  \begin{subfigure}[b]{0.48\textwidth}
    \centering
    \includegraphics[width=\textwidth]{figs/solucion_t1.pdf}
    \caption{$t = 0.5$: difusion parcial.}
    \label{fig:calor-t1}
  \end{subfigure}
  \caption{Evolucion de la ecuacion de calor
           $\eqop{\partial_t}\equ{u} = \eqk{\alpha}\eqop{\nabla^2}\equ{u}$
           con condiciones de Dirichlet.}
  \label{fig:calor-evolucion}
\end{figure}
\end{lstlisting}

La referencia a una subfigura usa el label de la subfigura:
\texttt{\textbackslash cref\{fig:calor-t0\}} produce \textit{fig.~7.1a}.
La referencia al grupo usa el label de la figura principal.

\subsection{Figuras embebidas en texto}
\label{subsec:wrapfig}

Para figuras pequeñas que deben fluir con el texto sin interrumpir el
párrafo, \texttt{wrapfig}~\cite{wrapfig2021} provee \texttt{wrapfigure}:

\begin{lstlisting}[style=latex, frame=none, numbers=none]
\begin{wrapfigure}{r}{0.35\textwidth}
  \centering
  \includegraphics[width=0.33\textwidth]{figs/diagrama.pdf}
  \caption{Diagrama del sistema.}
  \label{fig:diagrama}
\end{wrapfigure}
El texto circundante fluye alrededor de la figura. La opcion {r} la
coloca a la derecha; {l} a la izquierda. El ancho (0.35\textwidth)
debe ser ligeramente mayor que el de la imagen para dejar margen.
\end{lstlisting}

\begin{observacion}{Wrapfigure y listas}
\texttt{wrapfigure} no funciona correctamente dentro de entornos de lista
(\texttt{itemize}, \texttt{enumerate}) ni dentro de \texttt{tcolorbox}.
Tampoco debe aparecer al inicio de un capítulo antes de que exista
suficiente texto para fluir. Para esos contextos, usar \texttt{figure}
estándar.
\end{observacion}

%───────────────────────────────────────────────────────────────────────────────
\section{Gráficas con pgfplots}
\label{sec:pgfplots}
%───────────────────────────────────────────────────────────────────────────────

\begin{definicion}{Estilos Catppuccin para pgfplots}
\label{def:pgfplots-estilos}
\texttt{config/pgfplots.tex} define dos estilos con nombre y un ciclo
de colores de 8 series:
\begin{itemize}[leftmargin=2em]
  \item \texttt{catppuccin-clean} --- para publicación: grid mayor, fonts
    pequeñas, leyenda en esquina superior derecha.
  \item \texttt{catppuccin-filled} --- para notas o presentaciones: igual
    que \texttt{clean} pero con fondo \texttt{secondarybg}.
\end{itemize}
El ciclo de colores (\texttt{plot-series-1} a \texttt{plot-series-8})
garantiza que hasta 8 series sean visualmente distinguibles usando tanto
color como marcador distinto.
\end{definicion}

\subsection{Estructura básica de una gráfica}
\label{subsec:pgfplots-basico}

\begin{lstlisting}[style=latex, frame=none, numbers=none]
\begin{figure}[htbp]
  \centering
  \begin{tikzpicture}
    \begin{axis}[
      catppuccin-clean,
      xlabel = {$\eqk{t}$},
      ylabel = {$\equ{y}(\eqk{t})$},
      title  = {Decaimiento exponencial},
      xmin=0, xmax=5,
    ]
      \addplot[plot-series-1, thick] expression[domain=0:5]{exp(-x)};
      \addplot[plot-series-2, thick, dashed] expression[domain=0:5]{exp(-2*x)};
      \legend{$e^{-t}$, $e^{-2t}$}
    \end{axis}
  \end{tikzpicture}
  \caption{Comparacion de tasas de decaimiento.}
  \label{fig:decaimiento}
\end{figure}
\end{lstlisting}

\subsection{Gráfica con datos reales y barras de error}
\label{subsec:pgfplots-errorbar}

Para datos numéricos de simulaciones, \texttt{pgfplots} acepta coordenadas
directas con incertidumbre. El alias \texttt{plot-errorbar} controla el
color de las barras, que debe ser sutil para no competir con la serie:

\begin{lstlisting}[style=latex, frame=none, numbers=none]
\begin{axis}[
  catppuccin-clean,
  xlabel = {Numero de particulas $\eqk{N}$},
  ylabel = {$\langle \equ{E} \rangle / \eqk{N}$},
  xmode  = log,
]
  \addplot+[
    plot-series-1,
    error bars/.cd,
      y dir=both, y explicit,
      error bar style={plot-errorbar, thin},
  ] coordinates {
    (100,  -1.23) +- (0, 0.05)
    (500,  -1.41) +- (0, 0.02)
    (1000, -1.45) +- (0, 0.01)
    (5000, -1.47) +- (0, 0.005)
  };
\end{axis}
\end{lstlisting}

\subsection{Curva de nivel y mapa de calor}
\label{subsec:pgfplots-contour}

Los colormaps \texttt{catppuccin-heat} (magnitudes) y
\texttt{catppuccin-diverge} (residuos con signo) se definen en
\texttt{config/pgfplots.tex}:

\begin{lstlisting}[style=latex, frame=none, numbers=none]
% Mapa de calor de la funcion de densidad de probabilidad 2D
\begin{axis}[
  catppuccin-clean,
  colormap name  = catppuccin-heat,
  colorbar,
  view           = {0}{90},      % proyeccion superior
  xlabel         = {$\equ{x}_1$},
  ylabel         = {$\equ{x}_2$},
]
  \addplot3[surf, shader=interp, samples=40,
            domain=-3:3, domain y=-3:3]
    {exp(-0.5*(x^2 + y^2)) / (2*pi)};
\end{axis}

% Mapa divergente para residuos
\begin{axis}[
  colormap name = catppuccin-diverge,
  colorbar,
]
  \addplot3[surf, domain=-2:2] {x*y};
\end{axis}
\end{lstlisting}

\subsection{Gráfica 2D completa con theming semántico}
\label{subsec:pgfplots-completo}

Un ejemplo de gráfica funcional real con coloreado semántico en labels
y título. Muestra la convergencia del error de discretización
$\equ{e}_h = \|\equ{u}_h - \eqk{u}\|$ en función de $\eqk{h}$~\cite{tufte2001visual}:

\begin{lstlisting}[style=latex, frame=none, numbers=none]
\begin{figure}[htbp]
  \centering
  \begin{tikzpicture}
    \begin{loglogaxis}[
      catppuccin-clean,
      xlabel = {Paso de malla $\eqk{h}$},
      ylabel = {Error $\|\equ{u}_h - \eqk{u}\|_2$},
      title  = {Convergencia de diferencias finitas},
      legend pos = north west,
      xtick     = {0.0625, 0.125, 0.25, 0.5},
      xticklabels = {$1/16$, $1/8$, $1/4$, $1/2$},
    ]
      % Datos numericos: O(h^2)
      \addplot[plot-series-1, mark=*] coordinates {
        (0.5, 1.2e-2) (0.25, 3.1e-3) (0.125, 7.8e-4) (0.0625, 1.9e-4)
      };
      % Referencia teorica: pendiente -2 en log-log
      \addplot[plot-series-5, dashed, no markers]
        coordinates {(0.5, 1.2e-2) (0.0625, 1.9e-4)};
      \legend{FD centradas, $\mathcal{O}(\eqk{h}^2)$}
    \end{loglogaxis}
  \end{tikzpicture}
  \caption{Convergencia del esquema de diferencias finitas centradas
           para la ecuacion de Poisson 1D. La pendiente $-2$ en escala
           log-log confirma el orden teorico $\mathcal{O}(\eqk{h}^2)$.}
  \label{fig:convergencia-fd}
\end{figure}
\end{lstlisting}

\subsection{Subfiguras con pgfplots}
\label{subsec:pgfplots-subfig}

Para comparar resultados en múltiples paneles que comparten un tema
visual coherente:

\begin{lstlisting}[style=latex, frame=none, numbers=none]
\begin{figure}[htbp]
  \centering
  \begin{subfigure}[b]{0.48\textwidth}
    \begin{tikzpicture}
      \begin{axis}[catppuccin-clean, width=\textwidth, height=5cm,
                   title={Espacio de fase}, xlabel={$\equ{x}$},
                   ylabel={$\equ{\dot{x}}$}]
        \addplot[plot-series-1] table {data/fase.dat};
      \end{axis}
    \end{tikzpicture}
    \caption{Oscilador amortiguado: $\eqk{\gamma} < 2\eqk{\omega_0}$}
    \label{fig:fase-subamortiguado}
  \end{subfigure}
  \hfill
  \begin{subfigure}[b]{0.48\textwidth}
    \begin{tikzpicture}
      \begin{axis}[catppuccin-clean, width=\textwidth, height=5cm,
                   title={Serie temporal}, xlabel={$\eqk{t}$},
                   ylabel={$\equ{x}(\eqk{t})$}]
        \addplot[plot-series-2] table {data/tiempo.dat};
      \end{axis}
    \end{tikzpicture}
    \caption{Decaimiento amortiguado en el tiempo}
    \label{fig:tiempo-subamortiguado}
  \end{subfigure}
  \caption{Oscilador armónico subamortiguado: representación dual.}
  \label{fig:oscilador}
\end{figure}
\end{lstlisting}

%───────────────────────────────────────────────────────────────────────────────
\section{Captions: tipografía y colores}
\label{sec:captions}
%───────────────────────────────────────────────────────────────────────────────

El label del caption (``Figura 7.1'', ``Tabla 3.2'') usa el alias de color
del dominio \texttt{caption-*}, distinto por tipo de flotante. Esto permite
identificar visualmente de qué tipo es el elemento referenciado antes de
leer el número:

\begin{acadtable}[H]{tab:caption-colores}{Colores de labels de caption por tipo}
  \begin{tabularx}{\textwidth}{l l l X}
    \toprule
    \headerrow
    \thead{Tipo} & \thead{Alias} & \thead{Color Ctp} & \thead{Razón} \\
    \midrule
    Figura     & \texttt{caption-fig}  & \textcolor{caption-fig}{\textbf{Blue}}
      & = \texttt{link-internal}; coherente con referencias \\
    Tabla      & \texttt{caption-tab}  & \textcolor{caption-tab}{\textbf{Mauve}}
      & = \texttt{env-thm}; distinguible de figura \\
    Código     & \texttt{caption-code} & \textcolor{caption-code}{\textbf{Green}}
      & = \texttt{code-string}; evoca código \\
    Algoritmo  & \texttt{caption-alg}  & \textcolor{caption-alg}{\textbf{Peach}}
      & Cálido; distinguible de los anteriores \\
    Subfigura  & \texttt{caption-sub}  & \textcolor{caption-sub}{\textbf{Sapphire}}
      & = \texttt{env-phys}; un nivel por debajo de la figura \\
    \midrule
    \multicolumn{2}{l}{Texto descriptivo} & \textcolor{caption-text}{\textbf{Subtext0}}
      & Secundario; no compite con el label \\
    \bottomrule
  \end{tabularx}
\end{acadtable}

El texto descriptivo del caption (todo lo que sigue al label) usa
\texttt{caption-text} en gris suave para que el label coloreado ancle
visualmente el flotante sin que el texto completo compita con el contenido.

%───────────────────────────────────────────────────────────────────────────────
\section{linewidth vs textwidth}
\label{sec:linewidth}
%───────────────────────────────────────────────────────────────────────────────

\begin{definicion}{linewidth dentro de tcolorbox}
\label{def:linewidth}
Dentro de un entorno \texttt{tcolorbox}, \texttt{\textbackslash linewidth}
es el ancho del área de contenido de la caja (descontando los márgenes
\texttt{left=6pt} y \texttt{right=6pt} de \texttt{basebox}).
\texttt{\textbackslash textwidth} es siempre el ancho de la columna de texto
completa, independientemente del contexto. Usar \texttt{\textbackslash textwidth}
dentro de una caja produce overflow.
\end{definicion}

\begin{acadtable}[H]{tab:linewidth-contexto}{Ancho correcto según contexto}
  \begin{tabularx}{\textwidth}{l l X}
    \toprule
    \headerrow
    \thead{Contexto} & \thead{Ancho correcto} & \thead{Aplica a} \\
    \midrule
    Flujo normal    & \texttt{\textbackslash textwidth}
      & \texttt{figure}, \texttt{acadtable}, \texttt{tabularx} \\
    Dentro de tcolorbox & \texttt{\textbackslash linewidth}
      & \texttt{boxfig}, \texttt{boxtab}, \texttt{tabularx} \\
    Subfigura       & \texttt{\textbackslash textwidth} del subfigure
      & \texttt{\textbackslash includegraphics}, \texttt{tikzpicture} \\
    Wrapfigure      & ancho declarado en \texttt{wrapfigure}
      & imagen interna \\
    \bottomrule
  \end{tabularx}
\end{acadtable}

%───────────────────────────────────────────────────────────────────────────────
\section{PlaceFloat y FloatBarrier}
\label{sec:float-control}
%───────────────────────────────────────────────────────────────────────────────

El algoritmo de flotantes de \LaTeX\ puede acumular figuras y tablas no
colocadas hasta el final de una sección, dejando páginas desequilibradas.
Dos herramientas del template controlan esto:

\begin{lstlisting}[style=latex, frame=none, numbers=none]
% FloatBarrier: fuerza la colocacion de todos los flotantes pendientes
% antes de continuar. Util al inicio de una nueva seccion.
\FloatBarrier

% Placement [H] (paquete float): coloca el flotante exactamente aqui,
% sin algoritmo de flotantes. Usar cuando la figura es parte del argumento
% y debe aparecer en ese punto exacto del texto.
\begin{figure}[H]
  ...
\end{figure}
\end{lstlisting}

\begin{observacion}{Cuándo usar \texttt{[H]} y cuándo \texttt{[htbp]}}
\texttt{[H]} es para figuras que son parte del argumento narrativo ---
el lector necesita verlas en ese punto exacto para seguir el razonamiento.
\texttt{[htbp]} es para figuras que ilustran o complementan el argumento ---
el lector puede verlas antes o después sin perder el hilo. El uso excesivo
de \texttt{[H]} produce páginas con grandes espacios en blanco cuando la
figura no cabe en el espacio restante. \texttt{\textbackslash FloatBarrier}
al inicio de cada \texttt{\textbackslash section} es un buen balance.
\end{observacion}

%───────────────────────────────────────────────────────────────────────────────
\printbibliography[heading=subbibliography]
